\begin{otherlanguage}{french}

	\mtcselectlanguage{french}
	\setlocalecaption{french}{appendix}{Annexe}
	\setlocalecaption{french}{figure}{Figure}

	\chapter{Résumé en Français}
	\label{summary_in_french}

	\minitoc{}

	\section{Introduction}
	\label{sif_introduction}

	Les travaux présentés dans cette thèse s'inscrivent dans le contexte de la gestion des câbles pour
	la robotique sous-marine, et plus particulièrement dans le cadre du concept de \emph{cordée de
		robots} sous-marins. L'idée générale consiste à remplacer le schéma classique \enquote{un robot --
		un ombilical} par un système où plusieurs véhicules sont disposés le long du même câble et
	collaborent pour en contrôler la forme. Le robot de tête (leader) accomplit la mission principale
	(inspection, intervention, cartographie, etc.) tandis que les robots suiveurs régulent activement
	la géométrie du câble afin d'éviter le contact avec le fond, les obstacles et la coque du navire
	support.

	Du point de vue du contrôle, une telle cordée peut être vue soit comme un essaim de robots
	fortement contraints par un lien mécanique commun, soit comme un système articulé unique où chaque
	robot joue le rôle d'un maillon actionné relié aux autres par des liaisons flexibles et passives.
	Dans les deux cas, la géométrie du câble ne peut plus être traitée comme une perturbation
	secondaire : elle intervient directement dans la faisabilité des trajectoires, dans la sécurité
	(distance au fond, évitement d'emmêlement) et dans les charges appliquées aux véhicules.

	L'objectif central de cette thèse est donc de proposer des modèles de câble à la fois suffisamment
	précis pour représenter des ombilicaux et suffisamment rapides en termes de temps de calcul pour
	être utilisés en temps réel dans des lois de commande, puis de concevoir et d'évaluer des
	stratégies de contrôle de cordée qui exploitent explicitement ces modèles. Les contributions se
	structurent autour de deux axes complémentaires : la modélisation rapide de câbles pour la
	commande, notamment via une extension dynamique du modèle de chaînette, et la commande de chaînes
	de robots sous-marins reliés par des câbles, en tenant compte explicitement des contraintes
	géométriques et environnementales.

	Les travaux présentés dans cette thèse ont donné lieu aux publications et communications suivantes:

	\begin{itemize}
		\item \AtNextCite{\setcounter{maxnames}{99}}\fullcite{filliungAugmentedCatenaryModel2024}
		\item \AtNextCite{\setcounter{maxnames}{99}}\fullcite{peraudIMUbasedMonitoringBuoyBallast2024}
		\item \AtNextCite{\setcounter{maxnames}{99}}\fullcite{filliungAugmentedCatenaryModel2024a}
		\item \AtNextCite{\setcounter{maxnames}{99}}\fullcite{deowanSymbolicRegressionDynamic2025}
	\end{itemize}

	\section{État de l'Art}
	\label{sif_state_of_the_art}

	Dans un premier temps, il est nécessaire de replacer cette thèse dans le contexte plus général de
	la robotique sous-marine et des véhicules sous-marins sans pilote, largement présentée par exemple
	dans \textcite{creuzeRobotsMarinsSousmarins2014}. On distingue classiquement les \acp{rov}, reliés
	en temps réel à un navire support par un ombilical, et les \acp{auv}, qui opèrent de manière
	autonome avec des capacités de communication limitées. À côté de ces architectures
	conventionnelles, des véhicules plus spécialisés (biomimétiques, reconfigurables, hybrides)
	cherchent à améliorer la manoeuvrabilité, l'efficacité énergétique ou l'adaptabilité aux missions.
	Dans cette thèse, l'accent est mis sur les \acp{rov} en raison de la présence du câble, qui devient
	un élément central de la dynamique et du contrôle.

	\subsection{Robots sous-marins}
	\label{sif_sota_rov}

	Les \acp{rov} constituent aujourd'hui la plateforme privilégiée pour les missions de grande
	profondeur et les opérations industrielles, notamment dans les domaines de l'exploration et de
	l'inspection offshore \parencite{christROVManualUser2013}. Ils sont alimentés et pilotés à travers
	un ombilical reliant le robot au navire support, ombilical qui assure également une liaison de
	communication à haut débit. La nécessité de cette liaison filaire s'explique par la forte
	absorption des ondes électromagnétiques dans l'eau, qui rend les communications radio classiques
	inopérantes au- delà de quelques mètres
	\parencite{downingOpticalConstantsWater1975,rothmanHITRAN2004Molecular2005,altgeltWorldsLargestRadiostation2005}.

	Les \acp{rov} sont généralement classés en trois grandes familles : les engins de grande taille,
	capables de tâches complexes de construction ou de maintenance en grande profondeur ; les engins de
	taille intermédiaire, plus légers, opérables depuis des navires de taille modeste ; et les engins
	d'observation, ayant une masse réduite, principalement dédiés à l'inspection visuelle et à la
	surveillance environnementale
	\parencite{christROVManualUser2013,FlotteOceanographiqueFrancaise,EquipeRSMRobotique,StanfordRoboticsLab,khatibOceanOneRobotic2016}.

	\begin{figure}
		\centering
		\includegraphics[width=.48\figurewidth, trim=350 0 200 0, clip]{victor_6000.png}
		\includegraphics[width=.48\figurewidth, trim=150 0 170 0, clip]{atlante_and_victor_6000.png}
		\caption{
			Un \ac{rov} de grande taille: le Victor~6000 opéré par l'Ifremer.
		}
	\end{figure}

	\begin{figure}
		\centering
		\begin{subfigure}[b]{.48\figurewidth}
			\centering
			\includegraphics[width=\linewidth]{arthur.jpg}
			\caption{Un \ac{rov} de taille intermédiaire: Arthur, développé au LIRMM.}
		\end{subfigure}
		\begin{subfigure}[b]{.48\figurewidth}
			\centering
			\includegraphics[width=\linewidth]{betty.png}
			\caption{Un \ac{rov} d'observation: BlueROV Heavy ``Betty''.}
		\end{subfigure}
		\caption{}
	\end{figure}

	Les \acp{auv}, quant à eux, sont des plateformes autonomes qui ne sont pas reliées physiquement à
	la surface. Ils sont typiquement programmés en amont avec un plan de mission et s'appuient sur
	leurs capteurs embarqués, leurs algorithmes de contrôle et des communications acoustiques
	sporadiques ou des remontées à la surface pour accomplir leurs tâches (cartographie bathymétrique,
	suivi environnemental, recherche et sauvetage, etc.)
	\parencite{bechlioulisRobustFormationControl2019,merciSimulatorUnderwaterGlider2023}. Leur forme
	hydrodynamique, souvent de type torpille, est optimisée pour limiter la traînée et ainsi parcourir
	de longues distances \parencite{meyerGliderTechnologyOcean2016}.

	\begin{figure}
		\centering
		\includegraphics[width=\figurewidth]{asterix.jpg}
		\caption{
			Un \ac{auv} représentatif: AsterX opéré par l'Ifremer.
		}
	\end{figure}

	D'autres véhicules plus spécialisés existent, par exemple des robots biomimétiques s'inspirant de
	la nage des poissons, ou des systèmes reconfigurables capables de modifier la disposition de leurs
	propulseurs ou leur mode d'opération (hybride \ac{rov}/\ac{auv}) en fonction de la mission
	\parencite{triantafyllouEfficientSwimmingMachine1995,fagundesgasparotoAdvancesReconfigurableVectorial2021,raugelOperationalScientificCapabilities2019}.

	\begin{figure}
		\centering
		\includegraphics[width=\figurewidth, trim=250 100 250 100, clip]{ocean_one_k.png}
		\caption{
			Un \ac{rov} humanoïde de taille intermédiaire : Ocean One~K.
		}
	\end{figure}

	\subsection{Modélisation des câbles et systèmes couplés}
	\label{sif_sota_modelling}

	La modélisation des ombilicaux et des câbles joue un rôle central dans la commande des systèmes
	sous-marins encâblés. Le câble n'est pas seulement une contrainte mécanique : sa forme et les
	efforts qu'il transmet influencent directement la dynamique des véhicules, la faisabilité des
	trajectoires et la sécurité vis-à-vis du fond et des obstacles. Les approches de modélisation se
	positionnent sur un continuum allant de modèles géométriques simples à des modèles par éléments
	finis très détaillés.

	Les modèles géométriques, comme la chaînette (courbe hyperbolique) ou des paramétrisations
	ellipsoïdales, décrivent la forme du câble à partir de fonctions analytiques. Ils ont été largement
	validés en régime quasi-statique
	\parencite{druptValidityCatenaryModel2022,laranjeiraCatenarybasedVisualServoing2020} et présentent
	un coût de calcul très faible (de l'ordre de la milliseconde), ce qui les rend adaptés aux
	applications temps réel \parencite{smolentsevShapeVisualServoing2023}. En contrepartie, leur
	précision diminue dès que les effets dynamiques et les charges hydrodynamiques deviennent
	importants \parencite{filliungAugmentedCatenaryModel2024}.

	Une autre famille d'approches, dite de masse-ressort, discrétise le câble en une succession de
	masses reliées par des liaisons élastiques et amorties
	\parencite{hallValidationLumpedmassMooring2015,hongDynamicsModelingMotion2020}. Ce compromis entre
	fidélité et coût de calcul permet de capter des déformations importantes et des effets dynamiques,
	au prix d'un temps de calcul croissant avec le nombre de segments
	\parencite{dasilvagomesNewDynamicModel2021}.

	Les modèles par éléments finis (théorie de Kirchhoff, poutres géométriquement exactes, etc.)
	offrent la fidélité la plus élevée en prenant en compte la flexion, la torsion et de fortes
	déformations du câble
	\parencite{buckhamDevelopmentFiniteElement2004,chenDynamicCharacteristicsDeepsea2021,dehadraiTransientPlanarDynamics2022,beltranEvaluationCoupledEffect2020}.
	Ils sont adaptés à l'analyse fine et à la conception, mais leur coût de calcul (de la seconde à la
	minute) les rend difficilement utilisables dans des boucles de commande en temps réel.

	Des approches de couplage fin entre le véhicule et le câble s'appuient sur des modèles continus et
	des solveurs \ac{cfd} ou des méthodes \ac{sph}/\ac{dem} pour représenter les interactions
	fluide-structure de manière détaillée
	\parencite{wuIntegratedHydrodynamicsControl2018,xuHydrodynamicAnalysisTethered2022,xuNumericalSimulationHydrodynamics2022,suSPHDEMModelingCablecontrolled2024}.
	D'autres méthodes exploitent des formulations multi-corps avancées (\acp{ancf}, coordonnées mixtes)
	pour modéliser le système complet incluant le treuil, le câble et le véhicule, notamment lors des
	phases de déploiement et de récupération
	\parencite{htunSingularityfreeMultibodyModeling2021,htunTheoryApplicationAbsolute2022,banerjeeDeploymentControlCable1995,abreuMarineVehiclesStreamers2016}.
	Ces modèles fournissent une description précise des interactions, mais restent coûteux en calcul.

	Les travaux récents explorent aussi l'utilisation de \enquote{câbles intelligents} équipés de
	capteurs et d'actionneurs, capables par exemple de fournir des mesures de forme ou de modifier
	localement la flottabilité du câble afin de maîtriser sa géométrie
	\parencite{yuArmlessUnderwaterManipulation2013,vielROVLocalizationBased2023,vielSelfManagementROVUmbilical2022,vielEllipsoidTetherModel2024,KCFTechnologiesNaval}.

	\begin{figure}
		\centering
		\includegraphics[width=.5\figurewidth]{smart_tether.png}
		\caption{
			Exemple de câble instrumenté, équipé de capteurs inertiels.
		}
	\end{figure}

	\subsection{Contrôle des robots sous-marins}
	\label{sif_sota_control_uuv}

	Le contrôle des véhicules sous-marins autonomes et téléopérés est marqué par la présence de fortes
	non-linéarités, d'incertitudes sur les paramètres hydrodynamiques et de perturbations
	environnementales difficiles à modéliser. De nombreuses approches de commande robuste ont ainsi été
	proposées pour garantir la stabilité et la précision de suivi malgré ces incertitudes, comme le
	montrent par exemple des travaux combinant commande prédictive robuste et prise en compte des
	contraintes d'actionnement et d'obstacles \parencite{heshmati-alamdariRobustPredictiveControl2020}.

	Les méthodes de type \ac{nmpc} se sont imposées comme un cadre unificateur pour traiter les
	contraintes de positionnement, de vitesse et d'évitement d'obstacles dans des environnements
	complexes, tout en intégrant éventuellement des modèles simplifiés de la dynamique du véhicule. En
	parallèle, des approches hybrides mêlant contrôle en mode glissant et apprentissage statistique,
	par exemple via des processus gaussiens pour compenser les efforts hydrodynamiques mal modélisés
	\parencite{limaSlidingModeControl2020}, permettent d'améliorer le suivi de trajectoire en
	s'appuyant sur les données collectées en opération.

	Des techniques de commande adaptative basées sur des réseaux de neurones, parfois optimisées par
	des algorithmes évolutionnaires, ont également été proposées pour gérer les incertitudes
	paramétriques et les perturbations externes dans le contrôle de véhicules ou de manipulateurs sous
	marins \parencite{salloomAdaptiveNeuralNetwork2020}. À l'échelle de groupes de véhicules, des
	stratégies de contrôle de formation, fondées sur des fonctions de Lyapunov ou des cadres de
	commande distribuée, permettent de maintenir des formations d'\acp{auv} tout en suivant des
	trajectoires de référence \parencite{khalajiLyapunovBasedFormationControl2020}. Ces travaux
	s'accompagnent souvent de méthodes d'estimation d'état et de diagnostic de pannes robustes pour
	faire face au bruit de mesure, aux courants marins et aux dégradations d'actionneurs
	\parencite{songRobustStateEstimation2023}.

	\subsection{Contrôle des systèmes encâblés}
	\label{sif_sota_control_tethered}

	Les systèmes encâblés présentent des défis spécifiques liés au couplage dynamique entre le
	véhicule, le câble et l'environnement. Les premiers travaux sur les systèmes de gestion
	d'ombilicaux se sont concentrés sur des stratégies de déploiement et de récupération sûres, en
	mettant l'accent sur la fiabilité mécanique et la régulation de la tension dans le câble
	\parencite{abelUnderwaterVehicleTether1994,aaseModelingControlCatenary2018}.

	Au-delà de la seule robotique sous-marine, des études sur les systèmes aériens encâblés mettent en
	avant des architectures multi-robots où la longueur des câbles est activement régulée par des
	treuils et où le câble joue un rôle d'élément mécanique utile, par exemple pour le transport ou la
	manipulation d'objets
	\parencite{fagianoSystemsTetheredMulticopters2017,dantonioCatenaryRobotDesign2021}. Dans ces
	contextes, le système complet peut être vu comme un manipulateur articulé dont les segments sont
	des câbles de longueur variable et les véhicules des actionneurs aux extrémités.

	Dans le domaine spatial, les robots et satellites encâblés exploitent des cables pour des missions
	de transfert d'orbite, de désorbitation ou de manipulation de débris, avec des problèmes de
	commande spécifiques liés à l'absence d'amortissement atmosphérique et aux longues longueurs de
	câble \parencite{modiDeploymentDynamicsTether1979,huangDexterousTetheredSpace2017}. Des cadres de
	commande basés sur l'énergie et la passivité sont alors utilisés pour éviter l'emmêlement et
	réguler les oscillations du câble \parencite{wangAntitangleControlTethered2019}.

	Enfin, des résultats analytiques montrent que certains systèmes aériens encâblés sont plats au sens
	de la platitude différentielle, ce qui permet de concevoir des lois de commande non linéaires et
	des observateurs presque globalement convergents à partir d'ensembles de sorties judicieusement
	choisis \parencite{tognonDynamicsControlEstimation2017}. Ces travaux illustrent comment le câble
	peut être considéré non plus comme une simple perturbation, mais comme un élément de structure
	exploité explicitement dans la pose du problème de commande.

	\section{Modélisation d'un Câble}
	\label{sif_cable_modelling}

	\section{Contrôle de la Cordée de Robots}
	\label{sif_robot_chain_control}

	\section{Conclusion}
	\label{sif_conclusion}

\end{otherlanguage}